\section{Robust Median-of-Means Instantiation}

This section records Corollary S2.10, a robust finite-moment instantiation of the generic S2.8 confidence-envelope certificate.

Define
\[
Z_1=Y,
\qquad
Z_2=S,
\qquad
Z_3=YS,
\qquad
Z_4=S^2,
\qquad
Z_5=(U-Y)^2.
\]
Assume $S\ge0$, $0<\mathbb E[S]<\infty$, and
\[
0\le \operatorname{Var}(Z_j)\le v_j<\infty,
\qquad j=1,\ldots,5,
\]
where the $v_j$ are valid nonnegative population upper bounds.

Choose an odd integer $b$ satisfying
\[
b\ge8\log\frac{5}{\delta},
\qquad
b\le n,
\]
and set
\[
m=\left\lfloor\frac{n}{b}\right\rfloor.
\]
Use the first $bm$ held-out observations and split them into $b$ disjoint blocks of size $m$. For target $Z_j$, let $\bar Z_{j,r}$ be the mean in block $r$ and define
\[
\widetilde\mu_j
=
\operatorname{median}
\left(
\bar Z_{j,1},\ldots,\bar Z_{j,b}
\right).
\]
Set
\[
r_j=2\sqrt{\frac{v_j}{m}}.
\]

\begin{corollary}[Median-of-means S2.8 instantiation]
With probability at least $1-\delta$,
\[
|\widetilde\mu_j-\mathbb E[Z_j]|
\le r_j,
\qquad
j=1,\ldots,5.
\]
Define
\[
L_Y=\widetilde\mu_1-r_1,
\quad U_Y=\widetilde\mu_1+r_1,
\quad
L_S=\widetilde\mu_2-r_2,
\quad U_S=\widetilde\mu_2+r_2,
\]
\[
L_{YS}=\widetilde\mu_3-r_3,
\quad
U_{S^2}=\widetilde\mu_4+r_4,
\quad
U_M=\widetilde\mu_5+r_5.
\]
Let
\[
L_S^+=\max\{0,L_S\},
\quad
U_{S^2}^+=\max\{0,U_{S^2}\},
\quad
U_M^+=\max\{0,U_M\}.
\]
Let
\[
P_U=
\max\{L_YL_S^+,L_YU_S,U_YL_S^+,U_YU_S\},
\]
\[
C_L=L_{YS}-P_U,
\qquad
V_U=\max\{0,U_{S^2}^+-(L_S^+)^2\},
\]
and
\[
D_{\rm MoM}=C_L-\sqrt{U_M^+V_U}.
\]
Then
\[
P\!\left(
\operatorname{Cov}(U,S)\ge D_{\rm MoM}
\right)
\ge1-\delta.
\]
If $D_{\rm MoM}>0$ and $0<U_S<\infty$ on the simultaneous event, then
\[
\mathbb E_{\rm FP}[U]-\mathbb E[U]
\ge
\frac{D_{\rm MoM}}{U_S}>0.
\]
\end{corollary}

\begin{proof}
Fix target $Z_j$. A block mean satisfies
\[
\operatorname{Var}(\bar Z_{j,r})
\le\frac{v_j}{m}.
\]
If $v_j=0$, then nonnegativity of variance and the upper bound imply
\[
\operatorname{Var}(Z_j)=0.
\]
Hence $Z_j=\mathbb E[Z_j]$ almost surely, every block mean equals the population mean almost surely, $r_j=0$, and the MoM deviation event has probability zero.

Now suppose $v_j>0$. Chebyshev gives
\[
P\!\left(
|\bar Z_{j,r}-\mathbb E[Z_j]|>
2\sqrt{\frac{v_j}{m}}
\right)
\le\frac14.
\]
Call such a block bad. If the median-of-means estimator lies outside the interval of radius $r_j$, at least half the $b$ blocks must be bad. The block indicators are independent across disjoint i.i.d. blocks and have means at most $1/4$. Hoeffding therefore gives
\[
P\!\left(
|\widetilde\mu_j-\mathbb E[Z_j]|>r_j
\right)
\le
\exp\left(-\frac{b}{8}\right)
\le
\frac{\delta}{5}.
\]
The same displayed bound is trivially true in the $v_j=0$ case because its left-hand probability is zero. A union bound over the five target variables gives simultaneous coverage at least $1-\delta$. The resulting intervals are exactly the S2.8 inputs, so S2.8 yields
\[
\operatorname{Cov}(U,S)
\ge D_{\rm MoM}
\]
on the simultaneous event. The first-person lower bound follows from T1 and $\mathbb E[S]\le U_S$ when the certified covariance lower bound is positive.
\end{proof}

\paragraph{Moment requirement.}
This result requires finite variance of all five target variables. In particular,
\[
\operatorname{Var}(S^2)<\infty
\]
requires a finite fourth moment of $S$, while
\[
\operatorname{Var}((U-Y)^2)<\infty
\]
requires a finite fourth moment of the prediction residual. Likewise, $\operatorname{Var}(YS)<\infty$ requires $\mathbb E[Y^2S^2]<\infty$. The theorem therefore must not be summarized as assuming only finite variance of the raw variables.
